\subsection*{An explicit A.11 profile-size certificate from source stages}
\label{subsec:a11-explicit-profile-size-20261022}
The conditional whole-hull stress theorem previously required
$E_c\le\lambda/4$ as an independently supplied source enclosure.
For the exact Appendix-A stages, this requirement follows from
an explicit further smallness choice.

\begin{proposition}[Automatic smallness across both supports and their gap]
\label{prop:a11-source-E-small-20261022}
Assume the source (A.6)--(A.11), $P_*>e^{T_d}$,
$0<\lambda\le1/100$, the A.10 interpolation slope
$\ell\le-\lambda$, and the source-specific small A.11
root $|c_i|\le54\lambda^{28}$. Let $E_b$ be the uniform
precorrection A.11 plateau profile and
$E_c=E_b(1+\sum_i c_i\beta_i)$ for the two disjoint
width-$3/10$ source bumps with
$\|\beta_i\|_\infty\le640/3$.
Then everywhere on the entire closed A.11 two-support
hull (including the gap),
\[
0<E_c\le16P_*\lambda^{45}e^{-13/(2\lambda)}
\le16P_*\lambda^{45}.
\]
Consequently $64P_*\lambda^{44}\le1$ implies
$E_c\le\lambda/4$ uniformly.
\end{proposition}
\begin{proof}
At the inner reference $E\le P_*$. The one-unit initial
transition has $\ell\le3/5$, so multiplies $E$ by at most
$e^{1/10}$. The length-$T_d$ axial stage, one-unit
intermediate transition, and length-$60\log(1/\lambda)$
power-law stage give respectively $e^{-T_d/2}$,
$e^{-1/2}$ and at most $\lambda^{30}$.
Thus $E(X_p,\eta)\le P_*\lambda^{30}$.
The pulse has length $13/\lambda$ and slope $-\lambda$,
giving a factor at most $e^{-13/(2\lambda)}$.
The interpolation cannot increase $E$.
The first bump begins at plateau coordinate
$30\log(1/\lambda)-63/20$. Its baseline factor is at most
\[
e^{-(1/2+\lambda)(30\log(1/\lambda)-63/20)}
\le\lambda^{15}e^{693/400}<8\lambda^{15}.
\]
The last inequality uses $693/400<2$ and $e^2<8$,
with a rational Taylor-tail certificate.
The remaining baseline decreases across the whole hull.
Source bump disjointness gives
$|\sum_i c_i\beta_i|\le11520\lambda^{28}<1$.
Hence $E_c<2E_b$ and the theorem follows.
\end{proof}
\begin{corollary}[A cofinal parameter wedge]
\label{cor:a11-explicit-lambda-20261022}
If $T_f=512$ is permissible for the originating
A.10 profile and $P_*>e^{T_d}>1$, put
\[
Q=1+T_d+512+\log P_*,
\quad
\lambda=\frac{10^{-20}}{(1+P_*)Q},
\quad
h=\frac{\lambda e^{-T_d}}{10^6}.
\]
Then $0<\lambda<1/100$, $0<h<\lambda/100$, and
$64P_*\lambda^{44}<1$.
\end{corollary}
\begin{proof}
As $Q>1$, $\lambda<10^{-20}/(1+P_*)$ and
$64P_*\lambda^{44}
<64\cdot10^{-880}P_* /(1+P_*)^{44}<1$.
The other comparisons are immediate.
\end{proof}
This removes the separately assumed $E_c\le\lambda/4$
from prior conditional A.11 stress inequalities after
the authentic A.11 root bound and source stages are
instantiated. It does \emph{not} authenticate the same
profile's $S(\infty)=\Pi(\infty)=0$, future $D_s,P_s$,
complete exterior stress cone, Appendix B contraction,
unforced coefficient hierarchy, physical zero-force
residual or a trusted end-to-end Lean proof.
